% ============================================================================
% Chapter 29 --- Reversibility playbook
% ----------------------------------------------------------------------------
% Purpose : the institutional discipline that protects against one-way
%           commitments. Cross-cutting topic --- not a wave.
% Page target: ~4 pages.
% Dependencies: Ch. 17 (D17.6: reversibility non-negotiable);
%               Ch. 18-27 (each wave's reversibility criteria).
% ============================================================================

\chapter{Reversibility playbook}
\label{ch:reversibility}

The wave chapters of Part~III each declare their own reversibility
criteria. This chapter consolidates the cross-cutting reversibility
discipline: the reasoning behind it, the rollback procedures
applicable to each wave, the data and communication considerations
that any rollback engages, and the post-mortem framework that
captures what the institution has learned from a reversal.

%-----------------------------------------------------------------------------
\section{Why reversibility is non-negotiable}
\label{sec:rv-why}

The architectural commitment of this Reference is that no wave is
allowed to produce a one-way change at acceptable cost. The
commitment is not theoretical: every wave declares a reversibility
window, every component is procured with exit-assistance clauses
(D28.4), every change to the running state flows through the
declarative GitOps model (D14.4) so that prior states can be
reconstructed. The reasoning is structural rather than tactical.

\paragraph{Sovereignty as a temporal property.} Sovereignty is not
only the property of the architecture today; it is the property of
the architecture across time. An institution that has committed to
a transformation it cannot undo at acceptable cost has surrendered
its sovereignty over the matter, regardless of the transformation's
merits. The reversibility playbook is the operational expression of
this insight.

\paragraph{Reversibility differs from rollback.} A \emph{rollback}
is the execution of a documented procedure to return to a prior
state. \emph{Reversibility} is the standing capacity to do so within
an agreed cost envelope. The wave chapters declare reversibility
windows; this chapter develops the rollback procedures that those
windows guarantee.

\paragraph{When reversibility is forfeited.} The principle admits
exceptions when the institution's leadership accepts the loss of
reversibility explicitly, in writing, against the sovereignty grid
(\trchap{ch:sovereignty-grid}{}). The forfeit is documented with
the dimension on which it falls, the rationale, and a review
schedule. Forfeits are exceptional and are revisited at every
sovereignty grid annual review.

%-----------------------------------------------------------------------------
\section{Per-wave rollback procedures}
\label{sec:rv-procedures}

Table~\ref{tab:rv-procedures} summarises the rollback procedure
applicable to each wave. The procedures are not exhaustive; the
institution adapts them to its operational specifics. The intent is
to ensure that a rollback decision is not blocked by the absence of
a documented procedure.

\begin{table}[H]
\centering\small
\caption{Per-wave rollback procedures (indicative). The unit of
rollback is the smallest operational increment that the wave's
coexistence pattern admits.}
\label{tab:rv-procedures}
\begin{tabularx}{\textwidth}{%
  >{\hsize=0.5\hsize\linewidth=\hsize\bfseries\raggedright\arraybackslash}X%
  >{\hsize=0.6\hsize\linewidth=\hsize\raggedright\arraybackslash}X%
  >{\hsize=1.9\hsize\linewidth=\hsize\raggedright\arraybackslash}X%
}
\toprule
Wave & Rollback unit & Procedure summary \\
\midrule
0 — Foundations             & Mandate     & Governance body dissolved or reabsorbed; sovereignty grid de-adopted (informational state preserved). \\
1 — Identity                & Application cohort & Cohort returned to the incumbent identity provider; federation metadata refreshed; user re-onboarded if needed. Incumbent identity provider kept warm 6--12 months. \\
2 — Observability           & Collector source & New collectors removed at the source; legacy log paths unaffected throughout; parallel store decommissioned. \\
3 — Network                 & Segment / building & Segment configuration reverted to pre-cut-over version under change control; inter-segment routing rules updated. \\
4 — Policy                  & Policy rule & New bundle deployed reverting the rule to shadow or removing it; incumbent enforcement (conditional-access expression, in-app rule) re-instated. \\
5 — Pilot                   & Pilot     & Pilot decommissioned per Wave~5 stop procedure; data exported; tenancy released. \\
6 — Endpoint                & Population & Population returned to advisory mode; posture rule reverted in policy bundle. \\
7 — Audit and IR            & Runbook   & Refactored runbook reverted to prior version; refactored detections reverted in detection corpus. \\
8 — Workload classification & Workload  & Workload classification revised through labelling governance; cluster-separation schedule extended. \\
9 — Lifecycle automation    & Consumer group & Group returned to ticket-driven flow; self-service portal disabled for the group. \\
\bottomrule
\end{tabularx}
\end{table}

%-----------------------------------------------------------------------------
\section{Data preservation during rollback}
\label{sec:rv-data}

A rollback that succeeds technically but loses institutional data
has not actually succeeded. Three classes of data require explicit
preservation.

\paragraph{Operational data produced under the wave.} Audit
events, decision logs, configuration histories produced while the
wave was active are preserved according to the institution's
classification scheme. The rollback does not delete this data; it
only changes the operational state of the components that produced
it. The data remains available to subsequent investigation and to
the post-mortem of \S\ref{sec:rv-postmortem}.

\paragraph{User data affected by the wave.} Where a wave touches
user-facing storage or workspaces (Wave~5 pilot, Wave~8 workload),
the rollback preserves the user data through documented export
procedures. The export occurs before the rollback executes; the
exported data is provided to the user and retained according to the
institutional retention policy.

\paragraph{Migration artefacts.} Configuration that was authored
during the wave (policy-as-code rules, GitOps manifests, runbook
revisions) is not discarded by the rollback. It remains in version
control as a historical record and as a starting point for any
future re-attempt.

%-----------------------------------------------------------------------------
\section{Communication during rollback}
\label{sec:rv-communication}

Rollback communication addresses three audiences with three
distinct needs.

\paragraph{Affected user constituency.} The constituency directly
affected by the rollback (a faculty whose segment is reverted, a
research group whose pilot is decommissioned, a population whose
posture rule is reverted) receives advance notice through the
established communication channel for the wave. The notice
includes the rollback timeline, the operational impact, and the
named contact for questions.

\paragraph{Programme governance.} The governance body established
in Wave~0 is informed of the rollback at its next regular cycle,
or earlier if the rollback affects the programme's overall plan.
The information is decision-oriented: what triggered the
rollback, what was preserved, what the next steps are.

\paragraph{External stakeholders.} Where the rollback affects
external relationships (federation partners, supplier
contracts, regulatory bodies whose obligations are touched by
the wave), explicit communication ensures that the external
parties are not surprised. The message is calibrated to the
stakeholder: federation partners need technical detail,
suppliers need contractual clarity, regulators need
documentation that the institution remains within its
obligations.

%-----------------------------------------------------------------------------
\section{Post-mortem and re-planning}
\label{sec:rv-postmortem}

A rollback is followed by a structured post-mortem. The intent is
not to assign blame but to extract learning that the next attempt
can apply.

\paragraph{Post-mortem template.} The post-mortem records, at
minimum: the trigger that produced the rollback decision, the
rollback execution itself (what worked, what did not), the data and
communication outcomes, the root cause as best understood, and the
recommendations for any future attempt. The template is consistent
across waves so that post-mortems accumulate into a comparable
record.

\paragraph{Conditions for re-attempt.} The post-mortem produces
explicit conditions under which the wave (or its rolled-back
portion) may be re-attempted. The conditions typically include
specific technical changes, process changes, or pre-requisites that
were absent the first time. A re-attempt is treated as a new wave
phase, not as a continuation of the original.

\paragraph{Lessons-learned propagation.} The post-mortem is shared
with the broader programme, with peer institutions through the
\gls{nren} community, and where appropriate with the international
network anticipated in phase~2. The institutional cost of the
rollback becomes shared learning; the next institution that
attempts the same wave benefits from the experience.

\medskip

\noindent The reversibility playbook is the discipline that allows
the playbook of Part~III to be ambitious without being reckless.
An institution that has internalised the playbook can commit to a
wave with the confidence that, if the wave proves unworkable, the
institution will return to a known state at acceptable cost. That
confidence is what makes the architecture transformation a
governance decision rather than a leap of faith.
